% GENERATED FILE. Edit canonical lessons, then run npm run generate:books.
% canonical-source-hash: fnv1a64:09b7e688d35c13bf
% canonical-lessons: HI-C323-harana, HI-C323-pacharna, HI-C323-vijay, HI-C323-sahid, HI-C323-sandhi

\chapter{Victory and Defeat}
\label{ch:hi-victory-and-defeat}
\hlchaptermodality{323}

\begin{chapteropening}
\textbf{By the end of this chapter:} \emph{I can say who won a battle and how it ended.}

Complete the last lesson of chapter 323: I can say who won a battle and how it ended.
\end{chapteropening}

\section[\texorpdfstring{\hi{हराना} (harānā)}{हराना (harānā)}]{\hi{हराना} (harānā) — to defeat}

\label{lesson:HI-C323-harana}



\begin{quote}
Before the new one: say the Hindi for an arrow, then the Hindi for for the most part.
\end{quote}

\subsection*{You'll want to know: \hi{हराना}}
\textbf{\hi{हराना}} — \emph{harānā} — \textquotedblleft{}to defeat\textquotedblright{}.

\begin{quote}
\hi{छोटी} \hi{सेना} \hi{ने} \hi{बड़ी} \hi{सेना} \hi{को} \hi{हराया।}
\end{quote}

\begin{quote}
\emph{choṭī senā ne baṛī senā ko harāyā.}
\end{quote}

\begin{quote}
The small army defeated the big one.
\end{quote}

Past, with \hi{ने} (he or she …): \hi{उसने} \hi{दुश्मन} \hi{को} \hi{हराया}, \emph{usne duśman ko harāyā}.

\begin{cousinweb}
\textbf{\hi{हराना}} (\emph{harānā}) is the causative partner of \textbf{\hi{हारना}}, to lose, and CDIAL 14061 lists both under Sanskrit \emph{hārayati}: to make the other side lose.
\end{cousinweb}

\begin{grammarlens}[title={What you've built}]
One more word for telling what happened.
\end{grammarlens}

\subsection*{Your turn}
\begin{itemize}
\raggedright
  \item \emph{Say it:} \emph{harānā}
  \item \emph{Say it:} \emph{harānā}, once more
  \item \emph{Say it:} \emph{choṭī senā ne baṛī senā ko harāyā.} — The small army defeated the big one.
  \item \emph{Recall:} say the Hindi for an arrow, then the Hindi for for the most part, then say \emph{harānā} again
\end{itemize}

\begin{tcolorbox}[breakable,colback=teal!4,colframe=teal!35!black,title={Before you move on}]
What does \hi{हराना} mean? (\textquotedblleft{}To defeat\textquotedblright{}.) Say it once more.
\end{tcolorbox}
\section[\texorpdfstring{\hi{पछाड़ना} (pachāṛnā)}{पछाड़ना (pachāṛnā)}]{\hi{पछाड़ना} (pachāṛnā) — to throw down; to overpower}

\label{lesson:HI-C323-pacharna}



\begin{quote}
Before the new one: say the Hindi for for the most part, then the Hindi for to defeat.
\end{quote}

\subsection*{You'll want to know: \hi{पछाड़ना}}
\textbf{\hi{पछाड़ना}} — \emph{pachāṛnā} — \textquotedblleft{}to throw down; to overpower\textquotedblright{}.

\begin{quote}
\hi{एक} \hi{सिपाही} \hi{ने} \hi{दो} \hi{दुश्मनों} \hi{को} \hi{पछाड़ा।}
\end{quote}

\begin{quote}
\emph{ek sipāhī ne do duśmanoṁ ko pachāṛā.}
\end{quote}

\begin{quote}
One soldier overpowered two enemies.
\end{quote}

Past, with \hi{ने} (he or she …): \hi{उसने} \hi{दुश्मन} \hi{को} \hi{पछाड़ा}, \emph{usne duśman ko pachāṛā}.

\begin{cousinweb}
\textbf{\hi{पछाड़ना}} (\emph{pachāṛnā}) is inherited from \emph{*pracchaṭ}, a form CDIAL reconstructs (CDIAL 8493), with the sense of knocking down. It is the wrestler's word: to throw the other to the ground, and so to beat them.
\end{cousinweb}

\begin{grammarlens}[title={What you've built}]
One more word for telling what happened.
\end{grammarlens}

\subsection*{Your turn}
\begin{itemize}
\raggedright
  \item \emph{Say it:} \emph{pachāṛnā}
  \item \emph{Say it:} \emph{pachāṛnā}, once more
  \item \emph{Say it:} \emph{ek sipāhī ne do duśmanoṁ ko pachāṛā.} — One soldier overpowered two enemies.
  \item \emph{Recall:} say the Hindi for for the most part, then the Hindi for to defeat, then say \emph{pachāṛnā} again
\end{itemize}

\begin{tcolorbox}[breakable,colback=teal!4,colframe=teal!35!black,title={Before you move on}]
What does \hi{पछाड़ना} mean? (\textquotedblleft{}To throw down; to overpower\textquotedblright{}.) Say it once more.
\end{tcolorbox}
\section[\texorpdfstring{\hi{विजय} (vijay)}{विजय (vijay)}]{\hi{विजय} (vijay) — a victory}

\label{lesson:HI-C323-vijay}



\begin{quote}
Before the new one: say the Hindi for to defeat, then the Hindi for to throw down.
\end{quote}

\subsection*{You'll want to know: \hi{विजय}}
\textbf{\hi{विजय}} — \emph{vijay} — \textquotedblleft{}a victory\textquotedblright{}.

\begin{quote}
\hi{विजय} \hi{के} \hi{बाद} \hi{शहर} \hi{में} \hi{जुलूस} \hi{निकला।}
\end{quote}

\begin{quote}
\emph{vijay ke bād śahar meṁ julūs niklā.}
\end{quote}

\begin{quote}
After the victory a procession went through the city.
\end{quote}

Feminine.

\begin{cousinweb}
\textbf{\hi{विजय}} (\emph{vijay}) is Sanskrit \emph{vijaya}, \textquotedblleft{}victory\textquotedblright{} (Monier-Williams), the formal word beside \textbf{\hi{जीत}}. Hindi makes it feminine.
\end{cousinweb}

\begin{grammarlens}[title={What you've built}]
One more word for telling what happened.
\end{grammarlens}

\subsection*{Your turn}
\begin{itemize}
\raggedright
  \item \emph{Say it:} \emph{vijay}
  \item \emph{Say it:} \emph{vijay}, once more
  \item \emph{Say it:} \emph{vijay ke bād śahar meṁ julūs niklā.} — After the victory a procession went through the city.
  \item \emph{Recall:} say the Hindi for to defeat, then the Hindi for to throw down, then say \emph{vijay} again
\end{itemize}

\begin{tcolorbox}[breakable,colback=teal!4,colframe=teal!35!black,title={Before you move on}]
What does \hi{विजय} mean? (\textquotedblleft{}A victory\textquotedblright{}.) Say it once more.
\end{tcolorbox}
\section[\texorpdfstring{\hi{शहीद} (śahīd)}{शहीद (śahīd)}]{\hi{शहीद} (śahīd) — a martyr}

\label{lesson:HI-C323-sahid}



\begin{quote}
Before the new one: say the Hindi for to throw down, then the Hindi for a victory.
\end{quote}

\subsection*{You'll want to know: \hi{शहीद}}
\textbf{\hi{शहीद}} — \emph{śahīd} — \textquotedblleft{}a martyr\textquotedblright{}.

\begin{quote}
\hi{युद्ध} \hi{में} \hi{सौ} \hi{सिपाही} \hi{शहीद} \hi{हुए।}
\end{quote}

\begin{quote}
\emph{yuddh meṁ sau sipāhī śahīd hue.}
\end{quote}

\begin{quote}
A hundred soldiers fell in the war.
\end{quote}

Masculine or feminine, by the person meant.

\begin{cousinweb}
\textbf{\hi{शहीद}} (\emph{śahīd}) is usually traced, through Persian, to Arabic \emph{shahīd}, \textquotedblleft{}a witness\textquotedblright{} and then one who dies for a cause. \textbf{\hi{शहीद} \hi{होना}} is to fall for one's country.
\end{cousinweb}

\begin{grammarlens}[title={What you've built}]
One more word for telling what happened.
\end{grammarlens}

\subsection*{Your turn}
\begin{itemize}
\raggedright
  \item \emph{Say it:} \emph{śahīd}
  \item \emph{Say it:} \emph{śahīd}, once more
  \item \emph{Say it:} \emph{yuddh meṁ sau sipāhī śahīd hue.} — A hundred soldiers fell in the war.
  \item \emph{Recall:} say the Hindi for to throw down, then the Hindi for a victory, then say \emph{śahīd} again
\end{itemize}

\begin{tcolorbox}[breakable,colback=teal!4,colframe=teal!35!black,title={Before you move on}]
What does \hi{शहीद} mean? (\textquotedblleft{}A martyr\textquotedblright{}.) Say it once more.
\end{tcolorbox}
\section[\texorpdfstring{\hi{संधि} (sandhi)}{संधि (sandhi)}]{\hi{संधि} (sandhi) — a treaty; a peace pact}

\label{lesson:HI-C323-sandhi}



\begin{quote}
Before the new one: say the Hindi for a victory, then the Hindi for a martyr.
\end{quote}

\subsection*{You'll want to know: \hi{संधि}}
\textbf{\hi{संधि}} — \emph{sandhi} — \textquotedblleft{}a treaty; a peace pact\textquotedblright{}.

\begin{quote}
\hi{आख़िर} \hi{में} \hi{दोनों} \hi{राजाओं} \hi{ने} \hi{संधि} \hi{की।}
\end{quote}

\begin{quote}
\emph{ākhir meṁ donõ rājāoṁ ne sandhi kī.}
\end{quote}

\begin{quote}
In the end the two kings made a treaty.
\end{quote}

Feminine.

\begin{cousinweb}
\textbf{\hi{संधि}} (\emph{sandhi}) is Sanskrit \emph{saṁdhi}, \textquotedblleft{}a joining, an agreement\textquotedblright{} (Monier-Williams). Grammar uses the same word for the joining of sounds; history uses it for the joining of two sides.
\end{cousinweb}

\begin{grammarlens}[title={What you've built}]
One more word for telling what happened.
\end{grammarlens}

\subsection*{Your turn}
\begin{itemize}
\raggedright
  \item \emph{Say it:} \emph{sandhi}
  \item \emph{Say it:} \emph{sandhi}, once more
  \item \emph{Say it:} \emph{ākhir meṁ donõ rājāoṁ ne sandhi kī.} — In the end the two kings made a treaty.
  \item \emph{Recall:} say the Hindi for a victory, then the Hindi for a martyr, then say \emph{sandhi} again
\end{itemize}

\begin{tcolorbox}[breakable,colback=teal!4,colframe=teal!35!black,title={Before you move on}]
What does \hi{संधि} mean? (\textquotedblleft{}A treaty; a peace pact\textquotedblright{}.) Say it once more.
\end{tcolorbox}
